\documentclass{article}
\usepackage{azzam}

\begin{document}
\title{Weekly Self Practice OSK SMP Week 3}
\section*{Soal-Solusi}

\begin{enumerate}
    \item Tentukan digit ke $2019^{\text{th}}$ dari kiri pada bilangan $1223334444555556666667777777\dots$?
    
    \textbf{Penyelesaian.} Pertama, akan dicari angka berulang $n$ pada digit ke $2019^{\text{th}}$. Tinjau:
    \begin{align*}
        1+2+3+\dots+9+2(10+11+\dots+n) &< 2019 \\
        \frac{9 \times 10}{2} + 2\left(\frac{n(n+1)}{2} - \frac{9 \times 10}{2}\right) &< 2019 \\
        45 + 2\left(\frac{n^2+n}{2} - 45\right) &< 2019 \\
        45 + n^2 + n - 90 &< 2019 \\
        n^2 + n &< 2064 \\
        -45.934\dots < &n < 44.934\dots
    \end{align*}
    Sehingga bilangan asli terkecil $n$ yang mungkin adalah 44 dan diketahui $1+2+3+\dots+9+2(10+11+\dots+44)=1935$ maka di digit ke $1936^{\text{th}}$ angka 45 mulai berulang sebanyak 45 kali, tetapi selisih dari $2019-1935=84$ sehingga untuk hasil selisihnya ganjil akan menampilkan angka 4 dan untuk selisihnya genap akan menampilkan angka 5. Karena selisihnya 84 (genap) maka angka digit ke $2019^{\text{th}}$ nya adalah \fbox{5}.
    
    \item Berapa banyak cara memilih bilangan bulat berbeda $x, y$ dengan $0 \le x, y \le 19$ dan $5$ habis membagi $x+y$?
    
    \textbf{Penyelesaian.} Tinjau sisa pembagian setiap bilangan bulat dari 0 sampai 19. Ada 4 bilangan bulat di interval ini yang bersisa 0 jika dibagi 5; begitu pula dengan bilangan yang bersisa 1, 2, 3, dan 4 jika dibagi 5. Sekarang bagi 5 kasus berdasar sisa pembagian $x$ jika dibagi 5:
    \begin{itemize}
        \item Habis dibagi 5. Maka $y$ habis dibagi 5 pula untuk membuat $x+y$ habis dibagi 5. Perhatikan $x \ne y$: maka banyak cara memilih pasangan $(x, y)$ di kasus ini ada sebanyak $4 \times 3 = 12$.
        \item Bersisa 1 jika dibagi 5. Maka $y$ bersisa 4. Banyak cara: $4 \times 4 = 16$.
        \item Bersisa 2 jika dibagi 5. Maka $y$ bersisa 3. Banyak cara: $4 \times 4 = 16$.
        \item Bersisa 3 jika dibagi 5. Maka $y$ bersisa 2. Banyak cara: $4 \times 4 = 16$.
        \item Bersisa 4 jika dibagi 5. Maka $y$ bersisa 1. Banyak cara: $4 \times 4 = 16$.
    \end{itemize}
    Total ada sebanyak $12+16+16+16+16 = \fbox{76}$ cara.
    
    \item $x, y, z$ tiga bilangan yang memenuhi $(x+\frac{1}{y})(y+\frac{1}{z})(z+\frac{1}{x}) = 19$. Jumlah dari semua nilai yang mungkin dari $(x+\frac{1}{z})(z+\frac{1}{y})(y+\frac{1}{x})$ adalah \dots
    
    \textbf{Penyelesaian.} 
    \begin{align*}
        \left(y+\frac{1}{x}\right)\left(z+\frac{1}{y}\right)\left(x+\frac{1}{z}\right) &= \left(\frac{xy+1}{x}\right)\left(\frac{yz+1}{y}\right)\left(\frac{zx+1}{z}\right) \\
        &= \left(\frac{xy+1}{y}\right)\left(\frac{yz+1}{z}\right)\left(\frac{zx+1}{x}\right) \\
        &= \left(x+\frac{1}{y}\right)\left(y+\frac{1}{z}\right)\left(z+\frac{1}{x}\right) \\
        &= \fbox{19}
    \end{align*}
    
    \item Segitiga $ABC$ memiliki luas 120 satuan. Titik $E$ dan $F$ dipilih di sisi $AC$ sehingga $AE=EF=FC$. Jika $D$ dan $G$ berturut-turut merupakan titik tengah $AB$ dan $EF$, luas segitiga $DFG$ adalah \dots
    
    \textbf{Penyelesaian.} 
    $\frac{[DFG]}{[ABC]} = \frac{[DFG]}{[DAC]} \times \frac{[DAC]}{[ABC]} = \frac{FG}{AC} \times \frac{AD}{AB} = \frac{1/2 \times 1/3}{1} \times \frac{1/2}{1} = \frac{1}{12} \Rightarrow [DFG] = \frac{1}{12}[ABC] = \fbox{10}$.
    
    \item Nilai dari $\left(0,2^{\left(0,5^{\left(0,8^{\left(1,1^{\dots}\right)}\right)}\right)}\right)^{-1}$ adalah \dots
    
    \textbf{Penyelesaian.} Tinjau $0,8^{(1,1^{(1,4^{\dots})})}$ dan kita ketahui bahwa $\lim_{n \to \infty} x^n = 0$ untuk $|x|<1$ padahal $1,1^{1,4^{1,7^{\dots}}} \to \infty$ sehingga untuk $x=0,8$, maka $0,8^{1,1^{1,4^{\dots}}} \to 0$.
    Jadi, $\left(0,2^{\left(0,5^{\left(0,8^{\left(1,1^{\dots}\right)}\right)}\right)}\right)^{-1} \to (0,2^{(0,5^0)})^{-1} = (0,2^1)^{-1} = \left(\frac{1}{5}\right)^{-1} = \fbox{5}$.
    
    \item Bilangan $\frac{2019}{2^{2019}}$ merepresentasikan bilangan desimal. Tentukan digit ke-2017 di belakang koma dari bilangan pecahan tersebut.
    
    \textbf{Penyelesaian.} Nilai tersebut ekuivalen dengan mencari digit pertama dari 3 digit akhir dari $\frac{2019}{2^{2019}} \times 10^{2019} = 2019 \times 5^{2019} = 13 \times 31 \times 5^{2019}$. Akan dicari dengan modulo sehingga $13 \times 31 \times 5^{2019} \equiv x \pmod{1000}$. 
    Bagi menjadi 2 kasus:
    $x \equiv 13 \times 31 \times 5^{2019} \pmod{8} \equiv -3 \times (-1) \times 625^{504} \times 125 \pmod{8} \equiv 15 \times 1^{504} \equiv 7 \pmod{8}$.
    Juga $x \equiv 13 \times 31 \times 5^{2019} \pmod{125} \equiv 0 \pmod{125}$.
    Gabung kasus: karena kelipatan 125 maka $x=125a$. Tinjau $125a \equiv 7 \pmod{8} \Rightarrow 5a \equiv 15 \pmod{8} \Rightarrow a \equiv 3 \pmod{8} \Rightarrow a = 8b+3$.
    Jadi, $x = 125(8b+3) = 1000b + 375 \equiv 375 \pmod{1000}$. Digit pertama dari 3 digit akhir tersebut adalah \fbox{3}.
    
    \item Parabola $y=x^2$ dan $x=y^2$ berpotongan di titik $O=(0,0)$ dan $X=(1,1)$. Lingkaran $\omega$, yang berpusat di $(0,0)$ dan melewati $X$, berpotongan dengan parabola $y=x^2$ di titik $X$ dan $A$. Lingkaran $\omega$ berpotongan dengan parabola $x=y^2$ di titik $X$ dan $B$. Luas $\triangle XAB$ adalah \dots
    
    \textbf{Penyelesaian.} Pertama, jari-jari $\omega$ adalah $\sqrt{(1-0)^2+(1-0)^2} = \sqrt{2}$. Misal $A(p,q)$. Titik $A$ memenuhi persamaan $q=p^2$ dan $p^2+q^2=2 \iff p^2+p^4=2 \iff (p-1)(p+1)(p^2+2)=0$. Karena $A \ne X$, $p=-1 \Rightarrow A(-1,1)$.
    Dengan cara yang sama, didapat $B(1,-1)$. Perhatikan $\triangle XAB$ adalah segitiga yang siku-siku di $X$. Panjang $XA$ adalah 2 satuan, dan panjang $XB$ adalah 2 satuan. Luas $\triangle XAB$ adalah $\frac{2 \times 2}{2} = \fbox{2}$.
    
    \item Wadah pertama berisi 5 kelereng merah dan 4 kelereng biru. Wadah kedua berisi 7 kelereng merah dan 9 kelereng biru. Sebuah kelereng dipindahkan dari wadah pertama ke wadah kedua. Selanjutnya, sebuah kelereng diambil dari wadah kedua. Jika peluang yang terambil bola merah adalah $\frac{a}{b}$, tentukan nilai dari $a+b$.
    
    \textbf{Penyelesaian.} 
    Keadaan 1 (Pindah Merah): Peluang pindah merah = $\frac{5}{9}$. Peluang ambil merah dari wadah B = $\frac{7+1}{16+1} = \frac{8}{17}$. Peluang total Keadaan 1 = $\frac{5}{9} \times \frac{8}{17} = \frac{40}{153}$.\\
    Keadaan 2 (Pindah Biru): Peluang pindah biru = $\frac{4}{9}$. Peluang ambil merah dari wadah B = $\frac{7}{17}$. Peluang total Keadaan 2 = $\frac{4}{9} \times \frac{7}{17} = \frac{28}{153}$.
    Peluang yang dicari = $\frac{40}{153} + \frac{28}{153} = \frac{68}{153} = \frac{4}{9}$. Sehingga $a=4, b=9$. Jadi, $a+b = \fbox{13}$.
    
    \item $AT$ dan $BT$ menyinggung sebuah lingkaran berturut-turut di $A$ dan $B$, dan kedua garis bertemu tegak lurus di $T$. $Q$ dan $S$ titik berturut-turut di $AT$ dan $BT$, dan $R$ titik di lingkaran sehingga $QRST$ persegi panjang dengan $QT=8$, $ST=9$. Jari-jari lingkaran ini adalah \dots
    
    \textbf{Penyelesaian.} Perhatikan $TADB$ persegi, dengan $T$ titik pusat lingkaran. Misal $r$ radius lingkaran. Lihat segitiga $RID$. $|RI|=|AS|=|AT|-|ST|=r-9$, $|ID|=|QB|=|BT|-|QT|=r-8$, dan $DR=r$. Dari Phytagoras:
    $(r-9)^2+(r-8)^2 = r^2 \iff r^2-34r+145 = 0 \iff (r-29)(r-5) = 0$.
    Jelas $r>9$, maka yang memenuhi soal adalah $r=\fbox{29}$.
    
    \item Tentukan nilai dari:
    \begin{align*}
        \left(\frac{\sqrt{10+\sqrt{1}}+\sqrt{10+\sqrt{3}}+\sqrt{10+\sqrt{5}}+\dots+\sqrt{10+\sqrt{99}}}{\sqrt{10-\sqrt{1}}+\sqrt{10-\sqrt{3}}+\sqrt{10-\sqrt{5}}+\dots+\sqrt{10-\sqrt{99}}}\right)^2 +\\ \left(\frac{\sqrt{10-\sqrt{1}}+\sqrt{10-\sqrt{3}}+\sqrt{10-\sqrt{5}}+\dots+\sqrt{10-\sqrt{99}}}{\sqrt{10+\sqrt{1}}+\sqrt{10+\sqrt{3}}+\sqrt{10+\sqrt{5}}+\dots+\sqrt{10+\sqrt{99}}}\right)^2 
    \end{align*}
    
    \textbf{Penyelesaian.} Misalkan $A = \sum_{k=1}^{99} \sqrt{10+\sqrt{k}}$ dan $B = \sum_{k=1}^{99} \sqrt{10-\sqrt{k}}$. Tinjau $\frac{A}{B} = 1 + \frac{A-B}{B}$.
    $A-B = \sum_{k=1}^{99} \left(\sqrt{10+\sqrt{k}} - \sqrt{10-\sqrt{k}}\right) = \sum_{k=1}^{99} \sqrt{20-2\sqrt{100-k}} = \sqrt{2} \sum_{k=1}^{99} \sqrt{10-\sqrt{k}} = \sqrt{2}B$.
    Sehingga $\frac{A}{B} = 1+\sqrt{2}$. Jadi, soal $= \left(\frac{A}{B}\right)^2 + \left(\frac{B}{A}\right)^2 = (1+\sqrt{2})^2 + (\sqrt{2}-1)^2 = 3+2\sqrt{2}+3-2\sqrt{2} = \fbox{6}$.
\end{enumerate}

\end{document}